\chapter{Migraine}

\section*{Definition}
A primary headache disorder characterised by recurrent attacks of moderate to severe unilateral pulsating headache lasting 4-72 hours, often accompanied by nausea, photophobia, phonophobia, and sometimes aura.

\section*{What Happens in Your Body}
Migraine involves cortical spreading depression (CSD) activating the trigeminovascular system, releasing calcitonin gene-related peptide (CGRP) and substance P, causing neurogenic inflammation and widening of blood vessels of meningeal blood vessels. Oestrogen modulates trigeminal nerve sensitivity and CGRP release; the natural perimenopausal oestrogen decline and fluctuation can destabilise migraine patterns. The classic oestrogen withdrawal trigger occurs during the late luteal phase when oestrogen falls sharply. Perimenopause typically worsens migraine frequency in women with menstrual-related migraine, while post-menopause may improve it in some.

\section*{Causes}
\begin{itemize}
  \item Oestrogen fluctuation or withdrawal (perimenopause, menstrual cycle)
  \item Genetic predisposition (family history of migraine)
  \item Stress and anxiety
  \item Sleep disturbance (insufficient or excessive sleep)
  \item Dietary triggers (aged cheese, chocolate, alcohol, caffeine)
  \item Weather changes (barometric pressure, humidity)
  \item Hormonal medications (combined oral contraceptives, HRT)
  \item Dehydration or skipped meals
  \item Sensory triggers (bright lights, strong odours, loud noises)
  \item Medication overuse (analgesic rebound headache)
\end{itemize}

\section*{When to See a Doctor}
See your doctor if:
\begin{itemize}
  \item Symptoms are severe or getting worse over time
  \item Migraine with aura lasting more than 60 minutes, new-onset headache after age 50, worst headache of life, or headache with fever, stiff neck, or neurological deficits
  \item You have other concerning symptoms such as fever, weight loss, or bleeding
\end{itemize}

\section*{Related Symptoms}
\begin{itemize}
  \item Headache
  \item Nausea
  \item Vision blurred
  \item Photophobia
  \item Dizziness
\end{itemize}

\textbf{Important:} If this symptom persists, your doctor will check for serious underlying causes.

\section*{Self-Care / Home Management}
\begin{itemize}
  \item Lie down in a dark, quiet room at the first sign of a migraine and apply a cold pack to the forehead or neck.
  \item Identify and avoid personal trigger factors by keeping a headache diary for at least 1-2 months (track foods, sleep, stress, menstrual cycle, weather).
  \item Take acute medication at the earliest sign of an attack — triptans and NSAIDs are most effective when taken during the mild pain phase.
  \item Maintain a consistent sleep schedule (7-9 hours, same bedtime and wake time) as both too little and too much sleep can trigger migraines.
\end{itemize}

\section*{How Your Doctor Will Evaluate You}
\begin{itemize}
  \item Detailed headache history including onset, frequency, duration, quality, location, severity, aura features, triggers, and response to previous treatments.
  \item Neurological examination to identify focal deficits and exclude secondary headache causes (papilloedema, meningismus, cranial nerve abnormalities).
  \item No routine laboratory or imaging tests are needed for typical migraine with or without aura if the neurological exam is normal.
  \item Neuroimaging (CT or MRI brain) is indicated when there are atypical features: new-onset after age 50, thunderclap headache, progressive worsening, or abnormal neurological exam.
\end{itemize}

\section*{Treatment Options}
\begin{itemize}
  \item Acute attacks are treated with triptans (sumatriptan, rizatriptan, eletriptan) and NSAIDs; antiemetics (metoclopramide, domperidone) are used for nausea.
  \item Preventive therapy (propranolol, topiramate, amitriptyline, or candesartan) is indicated for frequent attacks (4+ per month) or attacks that significantly impair quality of life.
  \item CGRP monoclonal antibodies (erenumab, galcanezumab) or gepants (rimegepant, ubrogepant) are newer options for migraine prophylaxis and acute treatment respectively.
  \item Menopausal hormone therapy can stabilise migraine patterns in perimenopausal women with menstrual-related migraine; continuous combined HRT (rather than cyclical) is preferred to avoid oestrogen withdrawal.
  \item Medications, lifestyle modifications, and physical therapies may be prescribed as appropriate.
  \item Follow-up ensures treatment effectiveness and allows timely adjustments to the care plan.
\end{itemize}

\section*{References}

\begin{thebibliography}{99}
\bibitem{migraine:harrison-migraine} Longo DL, et al. \emph{Harrison's Principles of Internal Medicine}. 21st ed. McGraw-Hill; 2022. Chapter 425: Migraine.
\bibitem{migraine:uptodate-migraine} Silberstein SD, et al. Migraine in adults: acute treatment. In: UpToDate, Post TW (Ed), UpToDate, Waltham, MA; 2025.
\bibitem{migraine:headache-classification} Headache Classification Committee of the IHS. The International Classification of Headache Disorders, 3rd edition. \emph{Cephalalgia}. 2023;43(1):1-205.
\bibitem{migraine:dodick} Dodick DW, et al. Migraine: a clinical review. \emph{Lancet}. 2023;401(10380):860-874.
\bibitem{migraine:ashina} Ashina M, et al. CGRP and migraine: pathophysiology and therapeutics. \emph{N Engl J Med}. 2023;389(5):438-449.
\bibitem{migraine:aina} Ahn AH, et al. \emph{Wolff's Headache and Other Head Pain}. 9th ed. Oxford University Press; 2024. Chapter 11: Migraine.
\end{thebibliography}
